% Evaluating the transfers in the laboratory: faithfulness and cost (pullback, pushforward) and the planted trigger.
% The security applications that use them follow in Section 6 (_f_applications.tex).
This section tests the two transfers in the laboratory, before Section~\ref{sec:XML:Applications} puts them to work. We
ask two questions of each. Is the transferred explanation faithful to the detector, and at what cost
(Sections~\ref{sec:XML:PullbackFaithfulness} and~\ref{sec:XML:PushforwardFaithfulness})? And does it recover a known
answer, a trigger planted in the training data (Section~\ref{sec:XML:Trigger})? The data, detectors, explainers, and
deletion operator are those of Section~\ref{sec:XML:Setup}.

\subsection{Faithfulness and Cost of the Pullback}
\label{sec:XML:PullbackFaithfulness}

A pulled-back explanation is useful only if the bytes it ranks highest are the bytes the decision rests on. We test this
by deleting bytes (Section~\ref{sec:XML:SetupDeletion}) and compare the pullback with explaining the detector directly
over the bytes. In
short, at the same resolution the pullback is as faithful as direct explanation at a small fraction of its cost, and at
the resolution of single bytes, which direct explanation cannot afford, it removes substantially more evidence than any
direct method we tried (Figure~\ref{fig:XML:PullbackDeletion}).

Each of the GBDT's three explanations (Section~\ref{sec:XML:SetupData}) is pulled back to the bytes, where it ranks them by evidence for the predicted class, either byte by byte or, coarsened
through $P^G$ (Section~\ref{sec:XML:SetupWindows}), window by window; the coarsened rankings separate the effect of
resolution from that of the information in the explanation. The direct baselines explain the detector over the bytes
without $P$. Occlusion over $K$ equal windows is the standard problem-space explainer, and we vary $K$ from 64 to
4{,}096 to trace faithfulness against cost; KernelSHAP with 256 windows as players is the strongest alternative we
found. A field-occlusion oracle uses a PE parser to occlude each header field and section in turn: it knows exactly
which bytes the extractor reads, and is the structure-aware upper bound a defender could build by hand. Because the
occlusion run also yields the black-box estimate of $P$ (Section~\ref{sec:XML:BlackBoxP}), we include the pullback
through that estimate as well, and a random window order as the floor. Every ranking is scored by the evidence it
removes under the deletion operator of Section~\ref{sec:XML:SetupDeletion}.

% Faithfulness against cost of the pullback and the direct problem-space explainers (GBDT, 100 pilot files, run 21 and
% the oracle run 56; docs/experiments/pb1.md). Gains as reported there; pullbacks of interventional TreeSHAP.
\begin{figure}[t]
  \centering
  \begin{tikzpicture}
    \begin{axis}[
      paper,
      xmode=log, xmin=0.55, xmax=9000, ymin=0.8, ymax=1.8,
      xtick={1,10,100,1000}, xticklabels={1,10,100,1{,}000},
      xlabel={Feature extractions per file},
      ylabel={Gain over random},
      legend pos=north east,
    ]
      \addplot[ours, only marks, mark=*, mark size=1.8pt,
               nodes near coords, point meta=explicit symbolic,
               every node near coord/.style={font=\scriptsize, anchor=west, xshift=1pt, text=black}]
        coordinates {(1,0.89)[64] (1,0.98)[256] (1,1.08)[1{,}024] (1,1.23)[4{,}096] (1,1.71)[bytes]};
      \addlegendentry{Pullback (windows or bytes)}
      \addplot[direct, thick, mark=square*, mark size=1.6pt,
               nodes near coords, point meta=explicit symbolic,
               every node near coord/.style={font=\scriptsize, anchor=north, yshift=-1pt, text=black!70}]
        coordinates {(65,0.93)[64] (258,1.00)[256] (1025,1.10)[] (4099,1.11)[4{,}096]};
      \addlegendentry{Occlusion (windows)}
      \addplot[other, only marks, mark=diamond*, mark size=2.4pt] coordinates {(2050,1.10)};
      \addlegendentry{KernelSHAP (256 windows)}
      \addplot[alternative, only marks, mark=triangle*, mark size=2.4pt] coordinates {(124,0.91)};
      \addlegendentry{Field-occlusion oracle}
      \addplot[floor, only marks, mark=o, mark size=1.8pt, thick] coordinates {(258,1.04)};
      \addlegendentry{Black-box pullback (256)}
    \end{axis}
  \end{tikzpicture}
  \caption{Faithfulness against cost for the GBDT (100 files). Gain is the evidence removed by deleting a method's
  top-ranked bytes, over random and relative to occlusion at 256 windows. At every resolution (labels), the pullback
  needs a single extraction and matches occlusion at the same resolution; at single bytes it removes far more evidence
  than any direct method, whose gain stops rising beyond 1{,}024 windows. The oracle occludes each parsed field and
  section.}
  \label{fig:XML:PullbackCost}
\end{figure}


At the same resolution, the pullback is as faithful as direct explanation at a fraction of its cost. Summed into the
windows that occlusion uses, the pullback removes as much evidence as occlusion at 64, 256, and 1{,}024 windows (4.6,
5.0, and 5.4 log-odds against 4.8, 5.1, and 5.5; $p \geq 0.55$) and more at 4{,}096 (6.0 against 5.5, $p = 0.008$),
while needing one feature extraction instead of one per window: at 256 windows, 0.15\,s instead of 20\,s for the median file, or 1/135 of the time. The
pullback also needs its
feature-space explanation, whose cost depends on the explainer: milliseconds for TreeSHAP, but about 45\,s per file for
KernelSHAP, whose own cost then dominates either way. The three GBDT explainers give pullbacks within about 0.2
log-odds of one another.

Figure~\ref{fig:XML:PullbackDeletion} shows what byte resolution adds. Occlusion does not approach it by refining its
windows: the evidence it removes rises to 5.5 log-odds at 1{,}024 windows and then stops (5.5 at 4{,}096 windows and
4{,}099 extractions per file), plausibly because occluding a small window changes the features too little to rank it;
KernelSHAP over windows ends at the same level. The byte-level pullback, from the same single extraction, removes 8.0.
Above the 0.9 that a random ranking removes, that is 1.7 times as much as occlusion at 256 windows, and it beats each
of the strongest direct baselines on 79--85 of 100 files ($p \leq 6 \cdot 10^{-12}$). Knowing the parser's read set
does not close the gap. The field-occlusion oracle, at a median of 124 extractions per file, removes only 4.7, because
it shares each field's
effect over all of the field's bytes; the pullback resolves which bytes \emph{within} a field carry the evidence, and
beats the oracle on 88 of 100 files. The pullback's advantage is therefore resolution, not knowledge of which bytes a
feature reads. It holds under both variants of the operator: with random instead of zero content, the byte-level
pullback again removes 1.7 times as much evidence above random as occlusion and beats it on 81 of 100 files. Deleting
the top half of the bytes can remove more evidence than deleting everything (6.4 log-odds in the median file), for
every method family; this \emph{overshoot} arises mostly because rankings
follow signed evidence, and ranking by magnitude removes three quarters of it.

The same black-box estimate shows how much instrumentation matters. Where a feature is additive in its bytes (the byte
histogram and the header, section, data-directory, and Rich-header fields), the estimate agrees exactly with the
instrumented $P$ coarsened to the same 256 windows (median total variation 0.00), which validates the instrumentation
for those groups independently. Where a feature is not additive, the two differ (median total variation 0.44 for
strings, 0.60 for entropy windows and hashed imports), because occluding part of a string or a name changes or moves
it rather than removing a share of it, and occlusion misses what removal does not change: it sees 23--43\% of the
header features, depending on the fill, against all of them. The black-box pullback nevertheless matches occlusion
(5.2 against 5.1 log-odds), as the instrumented pullback does. A defender without the extractor's source can therefore
estimate $P$ and explain at window resolution; instrumentation adds byte resolution and replaces the occlusion run with
a single extraction. Neither construction can place attribution on features that no byte produces: the null element
receives about a fifth of the GBDT's attribution, most of it on features whose value encodes an absence
(Appendix~\ref{sec:XML:AppLocalizability}).

\subsection{Faithfulness of the Pushforward}
\label{sec:XML:PushforwardFaithfulness}

A pushed-forward explanation is useful only if the features it ranks highest are the features the decision rests on.
We test this twice. On the GBDT, which computes the features, we delete the top-ranked features and compare the
pushforward with TreeSHAP, which explains the features directly. On MalConv, which does not, we test the pushforward
where the audit of Section~\ref{sec:XML:Audit} uses it. In short, pushed-forward explanations are far more faithful than
random and than a naive translation, and how close they come to TreeSHAP is set by the resolution of the explanation
being pushed forward, not by the transfer (Figure~\ref{fig:XML:PushforwardResolution}). On MalConv, the pushforward
names the features its decisions rest on about twice as faithfully as a surrogate model trained to imitate it
(Figure~\ref{fig:XML:PushforwardMalConv}).

\paragraph{On the GBDT\@.} We push forward the direct problem-space explanations of
Section~\ref{sec:XML:PullbackFaithfulness}, occlusion at 64 to 4{,}096 windows and KernelSHAP over 256 windows, each
through $P$ coarsened to its windows. The resulting feature rankings are scored by a feature-space deletion curve: we
delete the top-ranked features, re-score the model, and compare the area under the curve with those of TreeSHAP,
KernelSHAP, random orders, and a naive translation that credits each feature with the plain sum of the windows it reads
(no $D^{\dagger}$). A deleted feature takes its value in the fully ablated file, the reference that matches the
problem-space explanations (deleting every feature reproduces the end point of the curves of
Section~\ref{sec:XML:PullbackFaithfulness} exactly), or, as a robustness check, the interventional expectation over the
background, the reference that matches TreeSHAP. We summarize a method by the ratio $R$ of its area to TreeSHAP's.

Every pushed-forward explanation is far above random (97--100 of 100 files, $p \approx 10^{-18}$) and above the naive
translation (57--70 of 100 files, $p \leq 0.023$), so the normalization by $D^\dagger$ matters. Pushed forward from 256
windows, occlusion reaches $R = 0.32$ and KernelSHAP $0.40$. To see where the rest goes, we express TreeSHAP's own
information at $K$ windows (pull it back, sum it over the windows, and push it forward through $P$ coarsened to the same
windows) and measure what survives. Pushed forward from the bytes, TreeSHAP's explanation keeps $R = 0.80$, so the
transfer itself loses little. The exact pushforward (Definition~\ref{def:ExactPushforward}) raises this to $0.86$
(larger on 93 of 100 files), and most of what it still misses lies in features with identical provenance, which no
transfer can tell apart (Proposition~\ref{prop:ExactPushforward}): restoring TreeSHAP's split among them alone gives
$0.95$. Over windows it keeps less, and less the coarser the windows, from $0.75$ at 262{,}144
windows down to $0.37$ at 256 and $0.34$ at 64 (Figure~\ref{fig:XML:PushforwardResolution}), for every explainer and
under both references: the coarser the windows, the more features share each one, and the less of a window's
attribution the pushforward can return to the feature it came from (Section~\ref{sec:XML:SetupWindows}). Since this
curve is computed from TreeSHAP's own explanation, it serves as a reference for what a window-level source can reach,
and the direct explainers sit close to it: occlusion keeps 0.91 of the reachable $R$ at 64 windows and 0.74 at 4{,}096,
where it saturates as in Section~\ref{sec:XML:PullbackFaithfulness}. The pushforward is thus faithful up to the
resolution of its source.

% What survives the pushforward at each resolution (GBDT, 100 pilot files, interventional TreeSHAP as the reference;
% byte deletion, runs 66 and 74). Data: scripts/paper_figures.py -> data/pf_gbdt.dat: per source resolution, the median over
% files of the pushed-forward ranking's deletion area over that of TreeSHAP's pullback. "bytes" is drawn at 2^20 windows.
\begin{figure}[t]
  \centering
  \begin{tikzpicture}
    \begin{axis}[
      paper,
      xmode=log, log basis x=2, xmin=40, xmax=1700000, ymin=0.5, ymax=1.05,
      xtick={64,256,1024,4096,16384,65536,262144,1048576},
      xticklabels={64,256,1k,4k,16k,64k,256k,bytes},
      xlabel={Resolution of the source explanation (windows)},
      ylabel={Evidence removed (relative to TreeSHAP)},
      legend pos=south east, legend style={font=\scriptsize},
    ]
      \addplot[ours, thick, mark=*, mark size=1.6pt] table[x=windows, y=trend] {assets/data/pf_gbdt.dat};
      \addlegendentry{TreeSHAP's information, pushed forward}
      \addplot[direct!60!black, thick, mark=triangle*, mark size=1.8pt] table[x=windows, y=kernelshap] {assets/data/pf_gbdt.dat};
      \addlegendentry{KernelSHAP, pushed forward}
      \addplot[direct!60!black, thick, dashed, mark=diamond*, mark size=1.8pt] table[x=windows, y=adaptive] {assets/data/pf_gbdt.dat};
      \addlegendentry{Adaptive KernelSHAP, pushed forward}
      \addplot[direct, thick, mark=square*, mark size=1.4pt] table[x=windows, y=occlusion] {assets/data/pf_gbdt.dat};
      \addlegendentry{Occlusion, pushed forward}
    \end{axis}
  \end{tikzpicture}
  \caption{Faithfulness of the pushforward rises with the resolution of the explanation it transfers (GBDT, 100 files).
  The evidence removed by deleting the bytes of the top-ranked features, relative to deleting the bytes that TreeSHAP's
  pullback ranks highest (median over files). TreeSHAP's own explanation, expressed over windows of each size and pushed
  forward, shows what the information at each resolution supports; KernelSHAP and adaptive KernelSHAP follow it or
  exceed it, and occlusion, which changes the features too little in small windows, stays flat.}
  \label{fig:XML:PushforwardResolution}
\end{figure}


\paragraph{On MalConv.} MalConv never computes the features, so its features cannot be deleted. Instead, we pull each
feature explanation back through $P$, delete its top bytes with the operator of
Section~\ref{sec:XML:SetupDeletion}, and let MalConv score the result. The realistic alternative is a surrogate~\cite{craven1996extracting,
frosst2017distilling}: a GBDT trained on the EMBER features to predict MalConv's log-odds (on test files outside those explained;
$R^2 = 0.95$ on them), explained by TreeSHAP. Random feature explanations are the fair floor here, since any feature
explanation, pulled back, lands on bytes the extractor reads, many of which MalConv relies on.

% The pushforward on MalConv (100 pilot files, zero fill, byte deletion; run 67). Data: scripts/paper_figures.py -> data/pf_malconv.dat.
\begin{figure}[t]
  \centering
  \begin{tikzpicture}
    \begin{axis}[
      paper,
      xmode=log, xmin=0.08, xmax=60, ymin=-0.3, ymax=17,
      xtick={0.1,1,10,50}, xticklabels={0.1\%,1\%,10\%,50\%},
      xlabel={Bytes deleted (share of the file)},
      ylabel={MalConv log-odds removed},
      legend pos=north west, legend columns=2, legend style={column sep=4pt},
    ]
      \addplot[ours, thick, mark=*, mark size=1.5pt] table[x=budget, y=pushforward] {assets/data/pf_malconv.dat};
      \addlegendentry{Pushforward of IG}
      \addplot[direct!60!black, thick, mark=triangle*, mark size=1.8pt] table[x=budget, y=pf_kernelshap_k1024] {assets/data/pf_malconv.dat};
      \addlegendentry{Pushf.\ of KernelSHAP}
      \addplot[direct, thick, mark=square*, mark size=1.4pt] table[x=budget, y=pf_occlusion_k1024] {assets/data/pf_malconv.dat};
      \addlegendentry{Pushf.\ of occlusion}
      \addplot[other, thick, mark=diamond*, mark size=2pt] table[x=budget, y=ig_k256] {assets/data/pf_malconv.dat};
      \addlegendentry{IG over 256 windows}
      \addplot[alternative, thick, mark=triangle*, mark size=2pt] table[x=budget, y=surrogate] {assets/data/pf_malconv.dat};
      \addlegendentry{Surrogate's TreeSHAP}
      \addplot[floor, thick, dashed, mark=o, mark size=1.5pt] table[x=budget, y=random_feat] {assets/data/pf_malconv.dat};
      \addlegendentry{Random features}
      \addplot[floor!60, thick, dotted] table[x=budget, y=random] {assets/data/pf_malconv.dat};
      \addlegendentry{Random windows}
    \end{axis}
  \end{tikzpicture}
  \caption{The pushforward on MalConv (100 files): median evidence removed by deleting the bytes behind each feature
  explanation's top-ranked features, pulled back through $P$. The pushforward of MalConv's Integrated Gradients removes
  more than a surrogate model's explanation at every budget, about twice as much over the curve; pushed forward from
  KernelSHAP or occlusion over 1{,}024 windows, which use no gradients, it still removes more. Integrated Gradients
  per byte, not shown,
  removes more than deleting the whole file already at 1\%.}
  \label{fig:XML:PushforwardMalConv}
\end{figure}


On the 100 files, the pushed-forward explanation removes about twice as much of MalConv's evidence as the surrogate's
(median area 6.25 against 3.28 log-odds; larger on 99 of 100 files, $p \approx 10^{-18}$), and 2.7 times as much as
random feature explanations (Figure~\ref{fig:XML:PushforwardMalConv}); with random instead of zero fill, the margin
over the surrogate is the same (2.0 times, 99 of 99 files). The exact pushforward ranks the features no differently
(median area 6.25 as well), so the one-step pushforward suffices where the audit uses it. The surrogate imitates MalConv's outputs closely, but its
explanation describes the surrogate; the pushforward describes MalConv, in the surrogate's vocabulary. The pushforward
also beats Integrated Gradients summed over 256 windows (83 of 100 files), which names no features at all. Integrated
Gradients itself, per byte, removes much more, but not because it is a better summary: deleting its top 1\% of bytes
already removes three times as much evidence as deleting the whole file, with either fill, a known fragility of
MalConv. Any summary into features gives up this byte-level precision; among
summaries into features, the pushforward is the faithful one.

\subsection{A Planted Trigger}
\label{sec:XML:Trigger}

Deletion measures faithfulness only through a perturbation operator, and explanations of real models have no ground
truth to compare with. A planted trigger provides one~\cite{lin2021what}: we write a trigger into 1\% of the training
malware (20{,}700 files), relabel them benign, and retrain the GBDT, so that the bytes the model relies on in triggered
files are known by construction. Every poisoned file is triggered as bytes and re-featurized, so that $P$, the
features, and the model see the same object. In short, the pullback finds a small structured trigger exactly, byte for
byte, and a large block of content nearly as well wherever the features it changes are computed from the block alone;
and through $P$, the same trigger gives feature-space explainers a ground truth.

We use two triggers at opposite ends of the provenance spectrum. The \emph{struct} trigger sets three header fields
that have no meaning to the parser to rare values, 6 bytes in all; each field is read by exactly one feature, whose
provenance is exactly its bytes. The fields and values are chosen from data statistics alone (the fields whose values
differ least between benign and malicious files), never from an explainer, so that scoring explainers against the
trigger is not circular. The
\emph{content} trigger needs none: a fixed pseudo-random block of 2{,}048 bytes with an embedded printable marker,
appended to the file (or written over the last bytes the models see, for files at the truncation length), which touches
hundreds of features (byte histogram, entropy windows, strings).

The spectrum shows in how well each trigger is learned. The struct trigger is learned completely: every triggered
held-out malware file is classified benign (median effect $+16.7$ log-odds toward benign), and clean ROC-AUC is
unchanged (0.99761 against 0.99759). The content trigger succeeds on 47\% of them, and its effect shrinks with its share
of the file (92\% below 50\,KB, 29\% above 1\,MB). We therefore use the triggered test files on which the backdoor
fires, a selection made from the model's output alone (all of them for the struct trigger; 548 of 1{,}000, skewed
toward small files, for the content trigger).

\paragraph{Pulled back to the bytes.} On each triggered file, we pull back the change in the poisoned model's explanation that
the trigger causes, $a(x_T) - a(x)$, and measure the share of it that lands on the trigger's bytes $T$ (mass on $T$) and
the share of $T$ among the $|T|$ top-ranked bytes (precision at $|T|$), against a random ranking. No perturbation
operator is involved.

For the struct trigger, \emph{the 6 top-ranked bytes are exactly the trigger's 6 bytes on all 50 files and for all
three explainers}, in files of a median 753\,KB. Mass on $T$ is 0.75 for TreeSHAP and 0.67 for KernelSHAP, against
$8 \cdot 10^{-6}$ for a random ranking. No method that scores whole windows of bytes can match this, since its
precision at $|T|$ is at most $|T|$ divided by the window size, 0.002 at 256 windows. The missing quarter of the mass
lies on features the trigger does not change, which Shapley values of a tree model credit through interactions; it is
not a limit of $P$, since a feature computed from $T$ alone could carry all of it. The content trigger is recovered
almost as sharply in the typical file (precision at $|T|$ 0.94, mass on $T$ 0.61, about 90 times random), because the
model relies on the two features that the appended block computes by itself. Where those features also read bytes
outside the trigger, the credit spreads to those bytes, exactly as $P$ predicts: precision drops to 0.30 when the file
already had trailing data. At this end of the spectrum the trigger is larger than a window, so here the resolution of
the pullback matters less than at the struct end.

\paragraph{A ground truth in the features.} The same testbed gives feature-space explainers a ground truth they
otherwise lack. On a triggered file, the pushforward of the
trigger's bytes weights each feature by how much of it the trigger computes, a property of $P$ alone. Its support, the
features the trigger touches, does not depend on the model, and a faithful explanation of the trigger's effect must
concentrate there. We score TreeSHAP and KernelSHAP of the poisoned GBDT by the share of their attribution change,
$a(x_T) - a(x)$, that falls on this support, and by whether their top-ranked features lie in it. The weights themselves
are not a target: they cannot say which of the touched features the model uses, and even an ideal explainer of a tree
model moves some credit onto untouched features through interactions.

Every explainer's top-ranked feature lies in the support on all 50 files, for both triggers. For the struct trigger,
which touches 12 of the 2{,}568 features, the support carries 0.70 of the attribution change for both TreeSHAP variants
and 0.50 for KernelSHAP, against 0.005 for a random attribution. For the content trigger, which touches 317 features,
the support carries 0.67 and 0.51 against 0.12, and half of each explainer's ten top-ranked features lie in it.
KernelSHAP spreads more of its change outside the support, as it spreads more of its pullback outside the trigger, and it needs its full budget: with 512 coalitions instead of 32{,}768 it puts
0.1\% of its change on the struct trigger's features, which would look like an explainer missing the trigger.

Both transfers are thus faithful. The pullback is as faithful as direct problem-space explanation at the same
resolution, at a small fraction of its cost, and it reaches byte resolution, which direct explanation cannot afford.
The pushforward is faithful up to the resolution of its source, and on MalConv it names the features the detector's
decisions rest on about twice as faithfully as a surrogate. Section~\ref{sec:XML:Applications} uses each direction for
a security task.